\documentclass[10pt,a4paper]{amsart}
\usepackage[margin=19mm]{geometry}
\usepackage{amsmath,amssymb,mathtools}
\usepackage[scr=rsfs,cal=euler]{mathalfa}
\usepackage{xfrac}
\usepackage[unicode,hidelinks,bookmarks=false]{hyperref}
\newcommand{\ie}{i.e.\ }
\newcommand{\cf}{cf.\ }
\newcommand{\resp}{respectively\ }
\newcommand{\Subsection}{\subsection}
\providecommand{\nobreakdash}{\nobreak\hbox{-}}
\setlength{\emergencystretch}{2em}
\multlinegap0pt
\usepackage{amssymb}
\usepackage{xcolor}
\usepackage{bm}
\newtheorem{lemma}{Lemma}
\title{On hyperbolic and functional analogues of questions of Gr\"unbaum and Loewner}
\author{Yu Huang, Sergii Myroshnychenko, Kateryna Tatarko, Vladyslav Yaskin}
\date{Source: arXiv:2606.04482v2 (CC BY 4.0)}
\begin{document}
\maketitle

\noindent\emph{Source and licence.} On hyperbolic and functional analogues of questions of Gr\"unbaum and Loewner. Version arXiv:2606.04482v2 (CC BY 4.0), distributed by arXiv under the Creative Commons Attribution 4.0 International licence, http://creativecommons.org/licenses/by/4.0/, as recorded in the arXiv licence metadata for this exact version and shown on the abstract page https://arxiv.org/abs/2606.04482v2. Full licence notice: \texttt{licenses/arxiv-2606.04482-CC-BY-4.0.txt}.\par\smallskip

\noindent\textit{Editorial scope.} Counted unit: the article's lemma Lemma:MR (source0.tex lines 480--485). The article's complete original proof of it is delivered with it (source0.tex lines 486--547, one \texttt{proof} environment of 62 lines). No mathematics is written, repaired or re-derived: the statement and the proof are the article's own lines, in the article's own notation, and no retained line is substituted or re-flowed.  Retained uncounted as the source-local context the proof needs: lines 260--262.  Nothing was omitted from any retained span, so the package carries no omission record.  No retained line contains a citation, so the wrapper carries no bibliography and no bibliography entry.  No retained line contains a figure environment, an image-inclusion command or drawing code, so no image is delivered and the package declares no asset dependency.  Editorial formatting only, in addition: an AMS \texttt{amsart} preamble that restates the article's own declarations and macros used by the retained lines, namely the environments \texttt{lemma} on separate counters, together with the standard packages those lines need.  The retained environments print in this document as Lemma 1 in order of appearance, whereas the article prints the same items with the numbering of its own full document.  The complete native source of the article as distributed in the arXiv e-print for this exact version is bundled unchanged as \texttt{sources/202051/source0.tex}.
\begin{equation}\label{s-conc}
	f(x)\le \frac{C}{1+|x|^{-1/s}},
\end{equation}

			\begin{lemma}\label{Lemma:MR}
				Let  $f\in C_s(\mathbb R^n)$ with $-1/(n+1)<s<\infty$  and $0 \in \text{int}\left(\text{supp}(f)\right)$. Then there exists at least one direction $u \in S^{n-1}$ such that 
				$$
				\int_{u^\perp} x f(x) \, dx = 0.
				$$
			\end{lemma}

			\begin{proof} Since an $s_1$-concave function is also $s_2$-concave for all $s_2<s_1$, we can assume that $f$ is   $s$-concave with $-1/(n+1)<s<0$.
				As was mentioned in the introduction, the integrals $$
				\int_{u^\perp} x f(x) \, dx 
				$$
				are well defined for all $u\in S^{n-1}$.
				
				Consider the following function of $u \in S^{n-1}$,
				$$
				F(u) = \int_{\{x: \, \langle x, u\rangle \geq 0\} } f(x) \, dx.
				%= \int_0^\infty \chi_{[0,\infty)}(x \cdot u) f(x) \, dx
				$$
				$F$ is continuous on $S^{n-1}$ since $f$ is integrable, and thus $F$ attains its extreme values. 
				
			Fix $u\in S^{n-1}$ and  let $v$ be a unit vector orthogonal to $u$. For a real number  $\varphi$ close to zero, define $$u_v(\varphi) = \cos\varphi\, u +\sin \varphi\, v.$$ We claim that 
				\begin{equation}\label{formula:der}
					\frac{d}{d\varphi} F(u_v(\varphi)) \Big|_{\varphi=0} =  \int_{u^\perp} \langle y,v\rangle f(y) \, dy.
					\end{equation}
				Indeed,
				\begin{align}\label{eqn:deriv}\frac{d}{d\varphi} F(u_v(\varphi)) \Big|_{\varphi=0} &=\lim_{\varphi\to 0} \frac1{\varphi} \left(F(u_v(\varphi))- F(u)\right)		\notag	\\
					&=\lim_{\varphi\to 0} \frac1{\varphi} \left( \int_{H^+} f(x) \, dx - \int_{H^-} f(x) \, dx\right)	,		
					\end{align}
						where $$H^+ = \{ x\in \mathbb R^n: \langle  x , u\rangle \leq 0 \ \text{and}\ \langle  x , u_v(\varphi)\rangle \geq 0 \}$$ and $$H^- = \{ x\in \mathbb R^n: \langle  x , u\rangle \geq 0 \ \text{and}\ \langle  x , u_v(\varphi)\rangle \leq 0 \}.$$
						It is enough to compute the limit in \eqref{eqn:deriv} as $\varphi\to 0^+$.  The case $\varphi\to 0^-$ will give the same result, since we can just replace $\varphi$ with $-\varphi$ and $v$ with $-v$.
						
						
						
						First, we will evaluate   $\lim_{\varphi \to 0^+}\frac{1}{\varphi}\int_{H^+} f(x) \, dx$.
							Note that any point $x \in H^+$ can be represented as $x = y + t u$, where $y \in u^\perp$, $ \langle y,v\rangle \ge 0$, and  $t \in [-\tan\varphi \langle y,v\rangle, 0]$.  Then
					\begin{align*}
					\lim_{\varphi\to 0^+} \frac1{\varphi}	\int_{H^+} f(x) \, dx & = \lim_{\varphi\to 0^+} \frac1{\varphi}\int\limits_{\{ y\in u^\perp:\langle y,v\rangle\ge 0\}} \left( \int\limits_{-\tan\varphi \langle y,v\rangle }^0 f(y + tu)\, dt \right) \,dy\\
					& = \lim_{\varphi\to 0^+} \frac1{\varphi}\int\limits_{\{ y\in u^\perp:\langle y,v\rangle\ge 0\}} \left(\int\limits_{-1 }^0 \tan\varphi \langle y,v\rangle f\left(y + \tan\varphi \langle y,v\rangle tu\right)\, dt \right)\, dy\\
					& = \int\limits_{\{ y\in u^\perp:\langle y,v\rangle\ge 0\}} \left( \int\limits_{-1 }^0 \lim_{\varphi\to 0^+}\left( \frac{\tan\varphi}{\varphi}  \langle y,v\rangle f\left(y + \tan\varphi \langle y,v\rangle tu\right)\right)\, dt \right)\, dy\\
					& = \int\limits_{\{ y\in u^\perp:\langle y,v\rangle\ge 0\}} \left(\int\limits_{-1 }^0   \langle y,v\rangle f\left(y \right) \, dt \right)\, dy\\
				& = \int\limits_{\{ y\in u^\perp:\langle y,v\rangle\ge 0\}}    \langle y,v\rangle f\left(y \right) \,   dy.
						\end{align*}						
							Above we used  the Dominated Convergence Theorem  to move the limit inside the integrals, since	by \eqref{s-conc} for small $\varphi$ 	we have					
				\begin{align*}
						\left| \frac{\tan\varphi}{\varphi}  \langle y,v\rangle f\left(y + \tan\varphi \langle y,v\rangle tu\right) \right| & \leq  2 |\langle y,v\rangle| f\left(y + \tan\varphi \langle y,v\rangle tu\right)\\
						& \leq  2 C |\langle y,v\rangle| {(1+|y + \tan\varphi \langle y,v\rangle tu|^{-1/s})^{-1}}\\
						& \le 2 C |\langle y,v\rangle| {(1+|y |^{-1/s})^{-1}},
					\end{align*}
						and the  latter is an integrable function on $u^\perp$.		
						
						
						The assumption that the origin lies in the interior of the support of 
						$f$  was used in computing the limit 	$$\lim_{\varphi\to 0^+}f\left(y + \tan\varphi \langle y,v\rangle tu\right)= f\left(y\right).$$ Since $f$ is continuous in the interior of its support, which is a convex set,   the equality above holds for almost every  $y\in u^\perp$, except possibly on a set of  zero $(n-1)$-dimensional measure.
						
						
						
					
						Thus, we obtain
						$$\lim_{\varphi\to 0^+} \frac1{\varphi}	\int_{H^+} f(x) \, dx = \int\limits_{\{ y\in u^\perp:\langle y,v\rangle\ge 0\}} \langle y,v\rangle f(y)\, dy.$$
						Similarly, one can show 
							$$\lim_{\varphi\to 0^+} \frac1{\varphi}	\int_{H^-} f(x) \, dx =  - \int\limits_{\{ y\in u^\perp:\langle y,v\rangle\le 0\}} \langle y,v\rangle f(y)\, dy.$$
						Formula \eqref{formula:der} follows by subtracting the last two limits, and recalling that the limit as $\varphi\to 0^-$ gives the same answer.
							
					 Let $u_0 \in S^{n-1}$ be a direction, where $F$ attains its minimum. Then,   formula \eqref{formula:der} implies
					 $$\int_{u_0^\perp} \langle y,v\rangle f(y) \, dy=0,$$ for any $v\in u_0^\perp$, which yields the statement of the lemma.
					 
			 
				
			\end{proof}

\end{document}
